\section{Growing radial support: simultaneous gain and its flux cancellation}
\label{cg:section}

This section carries a growing-support construction through the actual
radial Maxwell equations.  It proves simultaneous large gains for any
prescribed finite band of partial waves.  It also proves why this
particular construction does not give many power channels: the gains
of opposite electric inversion parities are reciprocal, whereas signed
flux through a centered section couples those opposite parities.
The resulting channel bound is independent of the mirror length and
radius within the prescribed band.  It is not a bound for arbitrary
dielectric geometry or for partial waves outside that band.

\subsection{A vacuum flux trace estimate with physical normalization}

\begin{lemma}[PROVED: trace-class vacuum section flux]
\label{cg:vacuum-trace}
Let $\mathcal H=L^2_t(S^2)$ be the power-normalized incoming vacuum
space.  For the regular vacuum Maxwell fields generated by
\eqref{hs:three-illumination}, let $Q_D^0$ be the Hermitian signed
flux operator through a bounded planar section of area $A$.  Then
\begin{equation}\label{cg:trace-bound}
 \|Q_D^0\|_1\leq\frac{2k_0^2A}{\pi}.
\end{equation}
If $D$ is centrally symmetric about the origin, the positive part has
\begin{equation}\label{cg:positive-trace}
 \operatorname{Tr}(Q_D^0)_+\leq\frac{k_0^2A}{\pi}.
\end{equation}
Consequently a subspace $V$ with
$\langle b,Q_D^0b\rangle\geq t\|b\|^2$ for all $b\in V$,
$t>0$, has dimension at most $k_0^2A/(\pi t)$ in the symmetric case.
\end{lemma}
\begin{proof}
Denote the two restriction maps to $L^2(D;\mathbb C^3)$ by
$\mathsf E$ and $\mathsf H$.  The physical normalization is
\begin{align}
 \mathsf E g(x)&=\frac{k_0\sqrt{Z_0}}{\sqrt2\,\pi}
                 \int_{S^2}g(d)e^{ik_0d\cdot x}\,dS(d),\notag\\
 \mathsf H g(x)&=\frac{k_0}{\sqrt{2Z_0}\,\pi}
                 \int_{S^2}d\times g(d)e^{ik_0d\cdot x}\,dS(d).
 \label{cg:herglotz}
\end{align}
The incoming far-field amplitude is
$\sqrt{2Z_0}\,g(-\widehat x)/r$ up to a unit-modulus phase, so
$\|g\|^2$ is incoming power.  At each direction the tangential
projection has trace two, and cross product by $d$ is an isometry
on its tangential range.  Direct integration of the squared kernels
therefore gives
\begin{equation}\label{cg:hilbert-schmidt}
 \|\mathsf E\|_{\rm HS}^2=\frac{4k_0^2Z_0A}{\pi},\qquad
 \|\mathsf H\|_{\rm HS}^2=\frac{4k_0^2A}{\pi Z_0}.
\end{equation}
For $Jv=v\times\widehat z$, $\|J\|\leq1$, and the flux operator is
\begin{equation}\label{cg:flux-factor}
 Q_D^0=\tfrac14
   (\mathsf E^*J\mathsf H+\mathsf H^*J^*\mathsf E).
\end{equation}
The product of two Hilbert--Schmidt operators is trace class, with
trace norm bounded by the product of their Hilbert--Schmidt norms.
The triangle inequality proves \eqref{cg:trace-bound}.

The electric inversion map $E(x)\mapsto E(-x)$ acts unitarily on
the incoming space.  Its two eigenspaces are the electric parity
sectors.  For a same-parity field the magnetic field has the opposite
parity by the curl equation, so its Poynting density is odd and its
flux through centered $D$ vanishes.  Polarization of this identity
makes both diagonal blocks of $Q_D^0$ zero.  The parity involution
therefore conjugates $Q_D^0$ to $-Q_D^0$.  Its positive and negative
spectra have equal multiplicities, and their trace norms are equal.
This proves \eqref{cg:positive-trace}.  The compact selfadjoint
min--max principle, followed by summing positive eigenvalues at least
$t$, proves the dimension assertion.
\end{proof}

\subsection{An explicit finite ordinary-dielectric stack}

Fix $1<q<\sqrt{\varepsilon_+}$, a positive integer $L$, and a
number $M\geq1$ of mirror periods.  The dimensionless radius is
$x=k_0r$.  Each period consists of a low-index interval of length
$\pi/2$ followed by a high-index interval of length $\pi/(2q)$.
The total dimensionless mirror length is
\begin{equation}\label{cg:length}
 d_M=\frac{M\pi}{2}(1+q^{-1}).
\end{equation}
Choose $x_0$ to be an integer multiple of $\pi$, and put
$x_1=x_0+d_M$.  The relative permittivity is vacuum for $x<x_0$
and $x>x_1$, and in a mirror interval with reference index
$n\in\{1,q\}$ it is
\begin{equation}\label{cg:material}
 \varepsilon(x)=n^2+2/x^2.
\end{equation}
This is a compact, real, scalar, nonmagnetic dielectric with finitely
many smooth radial pieces.  If
$x_0\geq\sqrt{2/(\varepsilon_+-q^2)}$, it obeys
$1\leq\varepsilon\leq\varepsilon_+$ everywhere.

Here are explicit finite-band radius requirements.  Write
\begin{equation}\label{cg:hankel-constants}
 p_l(x)=1+\sum_{s=1}^l a_{ls}x^{-s},\qquad
 S_l=\sum_{s=1}^l|a_{ls}|,\quad
 T_l=\sum_{s=1}^l s|a_{ls}|,\quad
 K_L=1+\max_{1\leq l\leq L}(2S_l+T_l),
\end{equation}
where $a_{ls}$ are the finite Hankel coefficients in
\eqref{ms:hankel-polynomials}, and set $D_L=L(L+1)+2$.
In addition to the material bound, require
\begin{equation}\label{cg:radius-choice}
 x_0\geq1,\qquad
 x_0\geq 8K_L q^{2M},\qquad
 x_0^2\geq32 D_Ld_Mq^{4M+4}.
\end{equation}
An integer multiple of $\pi$ satisfying these finitely many lower
bounds always exists.  The growing radius is an allowed resource
of this construction and is not hidden in a fixed-support claim.

\begin{theorem}[PROVED: simultaneous reciprocal gain of all prescribed partial waves]
\label{cg:simultaneous-gain}
For the dielectric \eqref{cg:material} with
\eqref{cg:radius-choice}, let $1\leq l\leq L$ and consider either
polarization.  Normalize its real regular radial solution in the
vacuum core to $u_l(x)=\psi_l(x)=xj_l(x)$, and write outside the
device
\begin{equation}\label{cg:outside-coefficient}
 u_l(x)=A_l\xi_l^-(x)+\overline{A_l}\xi_l^+(x).
\end{equation}
The coefficient is nonzero.  For TE use the electric radial
amplitude, and for TM the magnetic radial amplitude.  Uniformly
over the prescribed band and all azimuthal indices,
\begin{equation}\label{cg:amplitude-bounds}
 \frac{3}{16\sqrt2}\,q^{\sigma M}
   \leq |A_l|\leq\frac{5\sqrt2}{4}\,q^{\sigma M},
\end{equation}
where $\sigma=-1$ for TE odd $l$ and TM even $l$, and
$\sigma=+1$ for TE even $l$ and TM odd $l$.
Relative to the free-space regular solution with the same physical
incoming amplitude, the core multiplier is $g_l=(2A_l)^{-1}$.
Thus the electric-odd parity sector has gain $q^M$ times a complex
factor whose modulus lies between $2/(5\sqrt2)$ and $8\sqrt2/3$;
the electric-even sector has gain $q^{-M}$ times such a factor.
\end{theorem}
\begin{proof}
For TE take the dimensionless state $v=(u,u_x')^T$, and for TM
take $v=(u,\varepsilon^{-1}u_x')^T$.  These states are continuous
at every layer interface, and satisfy, respectively,
\begin{equation}\label{cg:generators}
 v'=\begin{pmatrix}0&1\\l(l+1)/x^2-\varepsilon&0\end{pmatrix}v,
 \qquad
 v'=\begin{pmatrix}0&\varepsilon\\l(l+1)/(\varepsilon x^2)-1&0\end{pmatrix}v.
\end{equation}
Their reference generators on an index-$n$ interval are obtained
by discarding the $x^{-2}$ terms and setting $\varepsilon=n^2$.
The exact quarter-period reference matrices are
\begin{equation}\label{cg:quarter-matrices}
 T_n^{\rm TE}=\begin{pmatrix}0&1/n\\-n&0\end{pmatrix},\qquad
 T_n^{\rm TM}=\begin{pmatrix}0&n\\-1/n&0\end{pmatrix}.
\end{equation}
The full reference products, in physical layer order, are
\begin{equation}\label{cg:reference-products}
 U_0^{\rm TE}(x_1,x_0)=(-1)^M\operatorname{diag}(q^{-M},q^M),
 \qquad
 U_0^{\rm TM}(x_1,x_0)=(-1)^M\operatorname{diag}(q^M,q^{-M}).
\end{equation}
Within an individual reference interval, each propagator and its
inverse has norm at most $q$.  Combining complete period products
with the at most two partial intervals gives the safe bounds
\begin{equation}\label{cg:prefix-bounds}
 \|U_0(x,x_0)\|,\ \|U_0(x,x_0)^{-1}\|\leq q^{M+2}
 \quad(x_0\leq x\leq x_1).
\end{equation}
For either polarization the difference of the actual and reference
generators has norm at most $D_L/x_0^2$.  For TE it is the single
entry $(l(l+1)-2)/x^2$.  For TM its entries are $2/x^2$ and
$l(l+1)/(\varepsilon x^2)$, with $\varepsilon\geq1$.

Writing the actual propagator as $U=U_0W$ gives a Volterra equation
whose coefficient norm integrates to at most
\begin{equation}\label{cg:beta}
 \beta=\frac{D_Ld_Mq^{2M+4}}{x_0^2}
       \leq\frac{q^{-2M}}{32}.
\end{equation}
Gronwall, followed by the exact endpoint norm in
\eqref{cg:reference-products}, gives
\begin{equation}\label{cg:transfer-error}
 \|U(x_1,x_0)-U_0(x_1,x_0)\|
       \leq q^M(e^\beta-1)\leq2q^M\beta.
\end{equation}
This estimate retains the possible growth of the reference transfer
matrix; no small generator error is identified with a small endpoint
error without this amplification factor.

The finite Hankel formulas give, for $x\geq1$,
\begin{equation}\label{cg:bessel-errors}
 |\psi_l(x)-\sin(x-l\pi/2)|
 +|\psi_l'(x)-\cos(x-l\pi/2)|\leq K_L/x.
\end{equation}
They also give $|\xi_l^+(x)|,|\xi_l^{+\prime}(x)|\leq2$
for $x\geq K_L$.  At the multiple $x_0$ of $\pi$, the ideal
initial state is a signed first coordinate vector for odd $l$ and
a signed second coordinate vector for even $l$.  Its error has norm
at most $K_L/x_0\leq q^{-2M}/8$, and the actual initial state
has norm at most two.  Combining this with
\eqref{cg:transfer-error} shows that the actual final state differs
from its ideal value in \eqref{cg:reference-products} by at most
\begin{equation}\label{cg:state-error}
 q^M\frac{K_L}{x_0}+4q^M\beta\leq\frac14q^{-M}.
\end{equation}
Therefore, with $\sigma$ as in the theorem,
\begin{equation}\label{cg:state-size}
 \tfrac34q^{\sigma M}\leq|v(x_1)|
       \leq\tfrac54q^{\sigma M}.
\end{equation}

In the exterior both state conventions equal $(u,u')$.  Since
$W(\xi_l^-,\xi_l^+)=2i$, the exterior coefficient satisfies
$A_l=(u\xi_l^{+\prime}-u'\xi_l^+)/(2i)$.
The Hankel bounds and the conjugate coefficient representation
give both inequalities
\begin{equation}\label{cg:coefficient-state}
 |v(x_1)|/(4\sqrt2)\leq|A_l|\leq\sqrt2|v(x_1)|.
\end{equation}
This proves \eqref{cg:amplitude-bounds}.  Nonvanishing also follows
directly from uniqueness: a zero incoming coefficient of a real
regular solution would make both exterior coefficients zero.

For a vacuum regular wave $\psi_l=(\xi_l^-+\xi_l^+)/2$, so its
incoming coefficient is $1/2$.  Fixed vacuum power normalization
therefore gives the exact relative core multiplier $1/(2A_l)$
for either polarization.  Finally the electric parity is $(-1)^l$
for TE and $(-1)^{l+1}$ for TM.  This identifies the amplified and
suppressed sectors as stated.
\end{proof}

\subsection{What happens to physical passage in this entire finite band}

\begin{theorem}[PROVED: an area bound survives the common-stack gains]
\label{cg:stack-count}
Let $D=(-L_D/2,L_D/2)^2\times\{0\}$, with area $A=L_D^2$, be
inside the vacuum core of the preceding construction.  Let
$\mathcal H_L$ contain all power-normalized incoming TE and TM
partial waves with $1\leq l\leq L$ and all their azimuthal indices.
For any physical selected-output contraction $T$ of this device,
every counted subspace $W\subset\mathcal H_L$ satisfies
\begin{equation}\label{cg:stack-dimension}
 \dim W\leq\frac{128}{9\pi}\,\frac{k_0^2A}{\eta\tau}.
\end{equation}
The bound is independent of $L,M,x_0$, the local resonant field norm,
and the exterior support size, provided the explicit radius conditions
are met.  Its proof uses no bypass hypothesis, so it applies in
particular whenever the frozen absolute-bypass condition also holds.
It does not assert a bound for inputs in angular orders above $L$.
\end{theorem}
\begin{proof}
Identify the finite physical incoming basis with the corresponding
orthonormal vacuum basis.  On the core, every actual field equals
the regular vacuum field with its coefficient in each mode multiplied
by $g_l$.  Let $P_-$ and $P_+$ be the electric-odd and electric-even
parity projections.  Theorem~\ref{cg:simultaneous-gain} writes this
diagonal multiplier exactly as
\begin{equation}\label{cg:gain-factor}
 G=(q^MP_-+q^{-M}P_+)B,
 \qquad \|B\|\leq C:=8\sqrt2/3,
\end{equation}
where every diagonal entry of $B$ is nonzero.  Its complex phases
are included in $B$; no common phase is assumed.

Both diagonal parity blocks of the vacuum section-flux operator
vanish, by the inversion argument in Lemma~\ref{cg:vacuum-trace}.
The two real gain factors in \eqref{cg:gain-factor} therefore cancel
in each cross block, giving the exact identity
\begin{equation}\label{cg:flux-cancellation}
 Q_D^{\rm device}\big|_{\mathcal H_L}
       =G^*Q_D^0G=B^*Q_D^0B.
\end{equation}
For counted inputs, $\langle a,Q_D^{\rm device}a\rangle
\geq\eta\tau\|a\|^2$.  Hence $V=BW$ has the same dimension as
$W$ and obeys
\begin{equation}
 \langle b,Q_D^0b\rangle
       \geq\frac{\eta\tau}{C^2}\|b\|^2\quad(b\in V).
\end{equation}
Apply \eqref{cg:positive-trace} and $C^2=128/9$.
This proves \eqref{cg:stack-dimension} for all coherent inputs,
without replacing the device transmission by an isolated link.
\end{proof}

\begin{theorem}[PROVED: arbitrary-rank physical field amplification with zero plane flux]
\label{cg:resource-counteraudit}
Fix $N\geq1$, a ball $B_R$ centered on the section plane, the same
frequency and material upper bound.  There is a sequence of the finite
ordinary dielectrics above and an $N$-dimensional physically
power-normalized input subspace $U_N$ such that the full scattering
map is unitary on $U_N$, every coherent field has
$S_z=0$ almost everywhere on the entire plane, and
\begin{equation}\label{cg:resource-growth}
 \int_{B_R}|E_a|^2\,dx
 \ \geq\ \frac{16Z_0}{25}q^{2M}
       \min_{1\leq j\leq N}
          \int_0^R|\psi_{2j-1}(k_0r)|^2\,dr\ \|a\|^2
 \qquad(a\in U_N).
\end{equation}
The positive constant after $q^{2M}$ is independent of $M$.
Thus every unit vector in this physical subspace has diverging local
field energy while full-map throughput is one and absolute bypass is
zero.  Its passage through $D$ is also zero, so it is not a counted
subspace for any $\eta>0$.
\end{theorem}
\begin{proof}
Take $L=2N-1$ and the TE modes $(l,m_{\mathrm{az}})=(2j-1,0)$,
$1\leq j\leq N$.  Choose $x_0$ also greater than $k_0R$ for each
$M$.  Radial real dielectrics preserve these orthogonal channels.
In each channel the outgoing-to-incoming coefficient is
$\overline{A_l}/A_l$, of modulus one, so their full physical
scattering map is unitary.

On the equatorial plane, $Y_{l0}=0$ for odd $l$ and
$\partial_\varphi Y_{l0}=0$.  A TE electric field there is purely
azimuthal, whereas its magnetic field is normal to the plane.
These component properties persist for every radius, in every layer,
and for every coherent sum of the selected channels.  Their normal
Poynting component is consequently zero everywhere except possible
interface circles, a set of plane measure zero.

A free unit-power TE partial wave has electric radial coefficient
$2\sqrt{2Z_0}\,\psi_l(k_0r)/r$ with unit-normalized $X_{l0}$.
The actual coefficient is multiplied by $g_l$.  The lower bound in
Theorem~\ref{cg:simultaneous-gain} gives
$|g_l|\geq2q^M/(5\sqrt2)$ for these odd orders.  Angular
orthogonality on $B_R$ removes all cross terms and proves
\eqref{cg:resource-growth}.  Each radial integral is strictly positive
because the regular Bessel solution is not identically zero.
The frozen passage inequality would require
$0=P_D(a)\geq\eta\|Ta\|^2>0$ for a nonzero efficient input;
thus no such input is counted.
\end{proof}

\begin{remark}[What the growing-support test establishes]
\label{cg:scope}
The simultaneous gain theorem closes a specific global synthesis step
for actual finite dielectrics: there is no assumed coupled-mode model,
no independent lumped loads, and no unproved extension of local Cauchy
data.  It also shows that an estimate of local field norm from
passivity, throughput, and bypass alone cannot be uniform, even on an
arbitrarily large fixed-dimensional physical subspace.

The positive-passage condition remains essential.  For this common
radial stack its required opposite-parity product cancels the large
gain, leading to the rigorous bandwise area bound
\eqref{cg:stack-dimension}.  To escape that result by angular-selective
inner layers one would have to change the parity-dependent radial
gain pattern while retaining full-plane bypass control.  The local
inverse-synthesis theorem does not prove this nonlocal statement.
Neither the construction nor its bandwise exclusion settles the
unrestricted channel conjecture.
\end{remark}
